\section{OPD Ablations}
\label{app:opd-ablations}

We ablate where the OPD loss is applied in the \qwenonesevenb{} OpenThoughts
comparison from Table~\ref{tab:qwen17_openthoughts_opsd_opd}. Vanilla OPD
applies the unprivileged teacher's loss to all sampled response tokens.
Epistemic-token OPD applies the same loss only to tokens in the
epistemic-marker set. Random-fraction OPD is a token-count-matched control: if
$x$ is the average fraction of epistemic tokens in student responses, then each
rollout receives OPD loss on a uniformly sampled $x\%$ subset of response
tokens. OPD + privileged gold-demonstration context uses the same loss over
response tokens but conditions the teacher on a gold demonstration.

\begin{table*}[!htbp]
\centering
\caption{\textbf{Token-masked OPD ablations do not reproduce the gold-demonstration degradation pattern.}
We evaluate \qwenonesevenb{} OPD variants on AIME24, AIME25, and HMMT25. Entries report pass@1 / pass@16 to match the full-format tables in Appendix~\ref{app:full-passk-results}. Epistemic-token OPD applies the loss only on epistemic-marker tokens. Random-fraction OPD applies the loss to a random fraction of response tokens matched to the average epistemic-token rate. Both token-masked OPD variants improve over the base model on average, while the gold-demonstration variant drops below the base on average at pass@1.}
\label{tab:opd_ablation_passk}
\footnotesize
\setlength{\tabcolsep}{3pt}
\renewcommand{\arraystretch}{1.12}

\newcommand{\metrichead}[1]{%
  \textbf{#1}%
}

\begin{tabular*}{\textwidth}{@{\extracolsep{\fill}}llcccc@{}}
\toprule
\textbf{Model}
&
& \metrichead{AIME24}
& \metrichead{AIME25}
& \metrichead{HMMT25}
& \metrichead{Average} \\
\midrule
\multirow{5}{*}{\textbf{\qwenonesevenb{}}}
& \textbf{Base}
& 0.502 / \textbf{0.800}
& 0.398 / 0.667
& 0.215 / 0.467
& 0.372 / 0.644 \\
&
\textbf{Vanilla OPD}
& \textbf{0.540} / \textbf{0.800}
& 0.385 / 0.667
& 0.252 / 0.567
& \textbf{0.392} / 0.678 \\
&
\textbf{Epistemic-only OPD}
& 0.523 / \textbf{0.800}
& 0.396 / 0.733
& 0.235 / 0.533
& 0.385 / 0.689 \\
&
\textbf{Random-frac OPD}
& 0.494 / \textbf{0.800}
& \textbf{0.408} / \textbf{0.767}
& \textbf{0.256} / 0.533
& 0.386 / \textbf{0.700} \\
&
\textbf{OPD + gold demo}
& 0.467 / \textbf{0.800}
& 0.344 / 0.667
& 0.240 / \textbf{0.600}
& 0.350 / 0.689 \\
\bottomrule
\end{tabular*}
\end{table*}


Table~\ref{tab:opd_ablation_passk} suggests that token-masked OPD can still
offer some improvement over the base thinking model even when the loss is
applied to only a small fraction of response tokens. The epistemic-token and
random matched-fraction masks are close enough that these results do not clearly
rank one mask above the other. This is consistent with the idea that lexical
epistemic markers such as \textit{wait} and \textit{hmm} are useful proxies for
forking behavior, but do not exhaust it: branch-relevant decisions can also
occur on ordinary mathematical, connective, or formatting tokens. A random
matched-fraction mask may therefore sample some consequential non-lexical
decision points, while the epistemic-token mask targets explicit deliberation
markers more directly.

\begin{table*}[t]
\centering
\caption{\textbf{Gold-demonstration context also lowers probability mass on epistemic tokens.}
This token-level companion to Table~\ref{tab:qwen17_openthoughts_epistemic_tokens} reports the model probability assigned to epistemic markers. The Marginal column gives aggregate probability mass on the epistemic-token set; the named columns give log-probabilities for representative markers; and Avg. logp averages over the set. Vanilla OPD leaves these probabilities nearly unchanged, while the gold-demonstration variant lowers both the aggregate marginal and several revision-token log-probabilities.}
\label{tab:qwen17_openthoughts_epistemic_logprobs}
\scriptsize
\setlength{\tabcolsep}{2pt}
\renewcommand{\arraystretch}{1.12}

\begin{tabular*}{\textwidth}{@{\extracolsep{\fill}}lccccccccc@{}}
\toprule
\textbf{Method}
& \textbf{Marginal}
& \textbf{wait}
& \textbf{recall}
& \textbf{okay}
& \textbf{altern}
& \textbf{check}
& \textbf{verify}
& \textbf{hmm}
& \textbf{Avg. logp} \\
\midrule
\textbf{Base}
& 0.00928
& -0.57
& -0.39
& -0.15
& -0.64
& -0.38
& -0.66
& -0.72
& -0.502 \\
\textbf{+OPD}
& 0.00929
& -0.57
& -0.38
& -0.15
& -0.64
& -0.38
& -0.67
& -0.70
& -0.499 \\
\textbf{+OPD gold demo}
& 0.00853
& -0.72
& -0.50
& -0.17
& -0.82
& -0.40
& -0.67
& -0.96
& -0.605 \\
\bottomrule
\end{tabular*}
\end{table*}


\begin{table*}[!htbp]
\centering
\caption{\textbf{Sparse-loss OPD controls leave epistemic-marker probabilities close to vanilla OPD.}
We report aggregate probability mass on the epistemic-marker set and log-probabilities for representative markers. Epistemic-token OPD and random-fraction OPD remain nearly identical to vanilla OPD on the aggregate marginal and average log-probability. Conditioning the teacher on a privileged gold demonstration lowers the marginal mass and assigns substantially lower probability to several revision markers, especially \textit{wait}, \textit{recall}, \textit{altern}, and \textit{hmm}.}
\label{tab:opd_ablation_epistemic_logprobs}
\scriptsize
\setlength{\tabcolsep}{2pt}
\renewcommand{\arraystretch}{1.12}

\begin{tabular*}{\textwidth}{@{\extracolsep{\fill}}lccccccccc@{}}
\toprule
\textbf{Method}
& \textbf{Marginal}
& \textbf{wait}
& \textbf{recall}
& \textbf{okay}
& \textbf{altern}
& \textbf{check}
& \textbf{verify}
& \textbf{hmm}
& \textbf{Avg. logp} \\
\midrule
\textbf{Base}
& 0.00928
& -0.57
& -0.39
& -0.15
& -0.64
& -0.38
& -0.66
& -0.72
& -0.502 \\
\textbf{Vanilla OPD}
& 0.00929
& -0.57
& -0.38
& -0.15
& -0.64
& -0.38
& -0.67
& -0.70
& -0.499 \\
\textbf{Epistemic-only OPD}
& 0.00928
& -0.57
& -0.38
& -0.15
& -0.64
& -0.38
& -0.66
& -0.71
& -0.498 \\
\textbf{Random-frac OPD}
& 0.00928
& -0.57
& -0.37
& -0.15
& -0.64
& -0.37
& -0.67
& -0.72
& -0.499 \\
\textbf{OPD + gold demo}
& 0.00853
& -0.72
& -0.50
& -0.17
& -0.82
& -0.40
& -0.67
& -0.96
& -0.605 \\
\bottomrule
\end{tabular*}
\end{table*}


\begin{table*}[!htbp]
\centering
\caption{\textbf{Gold-demonstration context produces the largest drop in realized epistemic-token density.}
The aggregate density column reports the fraction of generated tokens in the epistemic-marker set. The remaining columns report occurrences per 1,000 generated tokens for representative markers. Epistemic-token OPD and random-fraction OPD slightly reduce aggregate marker density relative to the base and vanilla OPD, but the gold-demonstration variant produces the largest decrease, including clear reductions in \textit{wait} and \textit{hmm}.}
\label{tab:opd_ablation_epistemic_tokens}
\scriptsize
\setlength{\tabcolsep}{2pt}
\renewcommand{\arraystretch}{1.12}

\begin{tabular*}{\textwidth}{@{\extracolsep{\fill}}lcccccccc@{}}
\toprule
\textbf{Method}
& \shortstack{\textbf{Epistemic token}\\[-1pt]\textbf{density}}
& \textbf{wait}
& \textbf{recall}
& \textbf{okay}
& \textbf{altern}
& \textbf{check}
& \textbf{verify}
& \textbf{hmm} \\
\midrule
\textbf{Base}
& 1.080\%
& 3.85
& 1.11
& 2.17
& 0.64
& 1.33
& 0.25
& 1.45 \\
\textbf{Vanilla OPD}
& 1.074\%
& 3.79
& 1.04
& 2.17
& 0.72
& 1.19
& 0.16
& 1.68 \\
\textbf{Epistemic-only OPD}
& 1.047\%
& 3.22
& 1.05
& 2.23
& 0.82
& 1.43
& 0.18
& 1.54 \\
\textbf{Random-frac OPD}
& 1.039\%
& 3.24
& 1.11
& 2.17
& 0.90
& 1.27
& 0.18
& 1.52 \\
\textbf{OPD + gold demo}
& 0.850\%
& 2.54
& 0.96
& 2.07
& 0.66
& 1.05
& 0.10
& 1.11 \\
\bottomrule
\end{tabular*}
\end{table*}

